\section{Introduction}
\label{sec:introduction}

Generalist visuomotor policies should reliably complete tasks across unseen
objects and complex states, adjusting subsequent actions based on observed
outcomes. Garment flattening and folding are well suited for evaluating these
capabilities. Garments vary widely in structure and appearance. Severe
crumpling causes self-occlusion, making the current configuration difficult to
infer. Completing the task requires both dynamic actions, such as flinging, and
precise local adjustments during flattening and folding.

Garment manipulation has been used to evaluate generalist policies on real
robots, but large-scale, reproducible comparisons remain difficult: data
collection, scene setup, and task resets require sustained human effort, and
garment initial configurations are difficult to reproduce precisely.
Simulation provides controlled experimental conditions and supports both data
reuse and scalable evaluation. However, existing simulation benchmarks provide
limited coverage of diverse garments and complex initial states, and their
datasets contain few successful trajectories that complete both flattening and
folding.

We present ClothDojo, a platform for data generation and policy evaluation in
garment flattening and folding. ClothDojo supports real-time dynamic
interaction between multiple bimanual robot embodiments and diverse garments.
Starting from severely crumpled states, closed-loop policies use flinging and
local adjustments to flatten garments, then continue to fold them. Unified
metrics measure task progress and final outcomes.

To scale the successful trajectories needed for downstream policy training,
ClothDojo decouples physical trajectory generation from appearance
randomization. A trajectory agent uses privileged visual information and is
initialized from human demonstrations. It is continually improved through
autonomous rollouts and human interventions. The resulting agent generates
successful closed-loop trajectories on new garments and initial states. We
then re-render these trajectories with randomized appearances while preserving
the underlying behavior, producing diverse RGB observations.

We use ClothDojo to study what enables generalist visuomotor policies to
complete this task. Paired experiments separate the effects of new physical
trajectories and appearance re-renderings, while controlled policy comparisons
test the roles of observation history, memory, and real-time control. We
evaluate all policies in closed loop on held-out initial states across seen and
unseen garments, reporting both stage-wise and end-to-end performance.
